% !TeX root = ../复分析笔记.tex
% 由主文档通过 \input 载入；请编译主文档。
\setcounter{chapter}{7} % 保留原书章号。
\chapter{洛朗级数与孤立奇点}
\label{chap:ca-laurent}

\section{为什么要引入负次幂}

泰勒级数适合研究中心点也属于定义域的函数。若函数只在 $0<|z-a|<R$ 全纯，中心可能出现 $1/(z-a)$ 一类发散项，普通幂级数就不够用了。允许负次幂后，可以同时描述中心的奇异行为和其余部分的正常变化。

\begin{theorem}[洛朗展开]
\label{thm:ca-laurent}
设 $f$ 在圆环 $A=\{r<|z-a|<R\}$ 全纯，其中 $0\le r<R\le\infty$，则有唯一展开
\begin{equation}
 f(z)=\sum_{n=-\infty}^{\infty}c_n(z-a)^n
 =\sum_{n=0}^{\infty}c_n(z-a)^n+\sum_{m=1}^{\infty}c_{-m}(z-a)^{-m}.
 \label{eq:ca-laurent}
\end{equation}
两部分级数在每个闭子圆环上绝对且一致收敛。对任意 $r<\rho<R$，
\begin{equation}
 c_n=\frac1{2\pi i}\oint_{|\zeta-a|=\rho}
 \frac{f(\zeta)}{(\zeta-a)^{n+1}}\,\mathrm d\zeta,
 \qquad n\in\mathbb Z,
 \label{eq:ca-laurent-coefficient}
\end{equation}
且系数与圆周半径 $\rho$ 无关。
\end{theorem}
\begin{proof}
取 $r<\rho_1<|z-a|<\rho_2<R$。对两圆之间的区域挖去 $z$ 附近的小圆，再用边界变形和柯西公式的证明，得
\[
 f(z)=\frac1{2\pi i}\oint_{|\zeta-a|=\rho_2}\frac{f(\zeta)}{\zeta-z}\,\mathrm d\zeta
 -\frac1{2\pi i}\oint_{|\zeta-a|=\rho_1}\frac{f(\zeta)}{\zeta-z}\,\mathrm d\zeta.
\]
外圆上的核按 $(z-a)/(\zeta-a)$ 展开，产生非负次幂；内圆上使用
\[
 \frac1{\zeta-z}=-\sum_{m=0}^\infty\frac{(\zeta-a)^m}{(z-a)^{m+1}},
\]
产生负次幂。对任何固定闭子圆环，把 $\rho_1,\rho_2$ 分别选在其内外两侧，两个几何级数都由严格小于 $1$ 的比值统一控制，故可逐项积分。系数积分的半径无关性由两圆之间的柯西定理得到。最后，若已有这种收敛展开，将其乘以 $(z-a)^{-n-1}$ 后沿圆逐项积分，式~\eqref{eq:ca-circle-power} 只保留第 $n$ 项，便得到唯一性和系数公式。
\end{proof}

式~\eqref{eq:ca-laurent} 中非负次幂部分称为正则部分，负次幂部分称为\emph{主部（principal part）}。双向求和的含义是这两个级数分别收敛后相加，不需要规定左右对称截断。实际计算通常不直接计算每个系数的积分，而是使用等比级数、已知泰勒级数、部分分式或有限次微分。

\section{同一函数在不同圆环有不同展开}

\begin{example}[展开式必须附带适用区域]
\label{ex:ca-annuli}
求 $f(z)=1/[(z-1)(z-2)]$ 以原点为中心的所有最大圆环展开。
\end{example}
\begin{solve}
先分解为 $f(z)=-1/(z-1)+1/(z-2)$。奇点距中心分别为 $1,2$，所以自然分成三个区域。在 $|z|<1$ 中，两项都按正幂展开：
\[
 f(z)=\sum_{n=0}^\infty\left(1-\frac1{2^{n+1}}\right)z^n.
\]
在 $1<|z|<2$ 中，第一项按 $1/z$ 展开，第二项按 $z/2$ 展开：
\[
 f(z)=-\sum_{n=0}^\infty z^{-n-1}
 -\sum_{n=0}^\infty\frac{z^n}{2^{n+1}}.
\]
在 $|z|>2$ 中，两项都按负幂展开：
\[
 f(z)=\sum_{n=0}^\infty(2^n-1)z^{-n-1}.
\]
最后一个展开的 $z^{-1}$ 系数为零，与 $f(z)$ 在无穷远处为 $O(z^{-2})$ 相符。中间圆环有无穷多负次幂，并不意味着原点是本性奇点：这个圆环不是以原点为唯一缺点的去心圆盘，不能用它的主部来分类原点。
\end{solve}

选择展开方式的原则是先确定哪一个比值的模小于 $1$。例如 $1/(z-1)$ 在 $|z|<1$ 应写成 $-1/(1-z)$，在 $|z|>1$ 应写成 $z^{-1}/(1-z^{-1})$。代数变形的选择来自区域条件，而不是来自形式喜好。

\section{孤立奇点的三种类型}

\begin{definition}[孤立奇点]
若 $f$ 在某个去心圆盘 $0<|z-a|<R$ 单值全纯，则称 $a$ 为其一个孤立奇点（isolated singularity）。按该去心圆盘上的洛朗展开，主部为零时称可去奇点；主部非零且只有有限项时称极点；主部有无穷多非零项时称本性奇点。
\end{definition}

\begin{theorem}[可去奇点判别]
\label{thm:ca-removable}
若 $f$ 在 $a$ 的去心邻域全纯，则以下条件等价：$a$ 是可去奇点；$f$ 能在 $a$ 补充一个值成为邻域内全纯函数；$f$ 在某个去心邻域有界；$\lim_{z\to a}f(z)$ 存在且有限。
\end{theorem}
\begin{proof}
主部为零时，非负次幂级数直接给出全纯延拓；延拓蕴含有限极限，有限极限蕴含局部有界。最后设 $|f|\le M$。由洛朗系数公式，对 $m\ge1$，
\[
 |c_{-m}|=\left|\frac1{2\pi i}\oint_{|z-a|=\rho}f(z)(z-a)^{m-1}\,\mathrm dz\right|
 \le M\rho^m.
\]
令 $\rho\downarrow0$ 得所有负次幂系数为零，完成等价性。
\end{proof}

例如 $\sin z/z$ 在原点可去，延拓值为 $1$；$(e^z-1)/z$ 也在原点可去，延拓值为 $1$。判断时不能仅看表达式是否“分母为零”，应看实际局部行为。

\begin{proposition}[极点判别]
\label{prop:ca-pole}
$a$ 为 $f$ 的 $m$ 阶极点，当且仅当
\[
 f(z)=\frac{g(z)}{(z-a)^m},\qquad g\text{ 在 }a\text{ 附近全纯且 }g(a)\ne0.
\]
此时 $|f(z)|\to\infty$ 当 $z\to a$，且 $1/f$ 可全纯延拓为在 $a$ 有 $m$ 阶零点的函数。反过来，若 $|f(z)|\to\infty$，则 $a$ 必为某个有限阶极点。
\end{proposition}
\begin{proof}
有限非零主部直接给出所列分解，分解也立即给出主部和倒数的零点。若 $|f|\to\infty$，则在足够小的去心邻域内 $f$ 无零点，$1/f\to0$，由可去奇点定理延拓。延拓后的倒数不恒为零，因此在 $a$ 的零点有有限阶，再取倒数即得极点。
\end{proof}

对有理函数 $P/Q$，若 $P$ 在 $a$ 的零点阶为 $p$，$Q$ 的零点阶为 $q$，则约分后由 $q-p$ 决定类型：$q>p$ 是 $q-p$ 阶极点，$q=p$ 可去且延拓值非零，$q<p$ 可去且延拓后为 $p-q$ 阶零点。计算极点阶数前，应先核对分子是否也为零。

\begin{example}[本性奇点的典型行为]
$e^{1/z}$ 在 $0<|z|<\infty$ 上展开为 $\sum_{n=0}^\infty z^{-n}/n!$，有无穷多负次幂，所以原点为本性奇点。沿正实轴趋于原点，模趋于无穷；沿负实轴趋于原点，函数趋于零；沿虚轴则持续振荡。它既不能补成有限值，也不像极点那样沿所有路径模都趋于无穷。
\end{example}

\begin{theorem}[本性奇点附近的稠密性]
\label{thm:ca-casorati}
若 $a$ 为 $f$ 的本性奇点，则在任意充分小的去心邻域中，$f$ 的像在 $\mathbb C$ 中稠密，即任意复数 $w$ 的任意小邻域都包含某个函数值。这称为卡索拉蒂--魏尔斯特拉斯定理（Casorati--Weierstrass theorem）。
\end{theorem}
\begin{proof}
若结论不成立，则存在 $w\in\mathbb C$、$\delta>0$ 及某个去心邻域，使 $|f(z)-w|\ge\delta$。于是 $g=1/(f-w)$ 在此全纯且有界，可延拓到 $a$。若 $g(a)\ne0$，则 $f=w+1/g$ 可去；若 $g(a)=0$，则 $g$ 不恒为零，零点阶有限，$f$ 在 $a$ 为极点。两种情况都与本性奇点矛盾。
\end{proof}

稠密不等于每一个值都能取到。例如 $e^{1/z}$ 永不取零，但可以无限接近零。更强的取值结论属于后续复分析课程，不是这里证明的内容。

\section{无穷远点与亚纯函数}

设 $f$ 在 $|z|>R$ 全纯，通过 $F(\zeta)=f(1/\zeta)$ 在 $0<|\zeta|<1/R$ 的行为定义 $f$ 在无穷远点的奇点类型。次数为 $n\ge1$ 的多项式在无穷远点有 $n$ 阶极点；$1/z$ 在无穷远点可去，延拓值为零；$e^z$ 在无穷远点为本性奇点。

一个函数在区域内除孤立极点外全纯，称为\emph{亚纯函数（meromorphic function）}。有理函数是亚纯函数的基本例子；$\tan z$ 在平面内亚纯，因为其奇点都是极点。亚纯不允许本性奇点出现在区域内部。

若整函数在无穷远点可去，则它在大圆外有界，在大圆内也有界，由刘维尔定理为常数；若在无穷远点有 $m$ 阶极点，则 $|f(z)|\le C|z|^m$ 对充分大的 $|z|$ 成立，柯西估计使高于 $m$ 阶的泰勒系数全部为零，故它是多项式。这个例子说明，研究无穷远点可以把函数的整体增长转化为一个局部奇点问题。

\begin{theorem}[球面上的亚纯函数]
\label{thm:ca-sphere-rational}
一个函数在整个黎曼球面上亚纯，当且仅当它是有理函数；这里不把恒等于无穷的映射称为亚纯函数。
\end{theorem}
\begin{proof}
有理函数在有限点只有极点或可去奇点，在无穷远处令 $z=1/\zeta$ 后仍是有理式，故在球面上亚纯。反之，球面上亚纯函数 $f$ 的极点没有聚点，否则聚点附近不能具有全纯或孤立极点的局部表示。由球面紧性，极点只有有限个。减去各有限极点的有限洛朗主部，得到整函数 $g$；这些主部在无穷远处都趋于零，所以 $g$ 在无穷远处仍可去或为极点。由上段结论，$g$ 是多项式。因此 $f$ 是一个多项式与有限个主部之和，恰为有理函数。
\end{proof}

\section{习题}
\label{sec:ca-laurent-ex}
\begin{enumerate}
 \item 分别在 $0<|z|<1$ 和 $|z|>1$ 展开 $1/[z(z-1)]$。
 \item 分类原点处的奇点：$(1-\cos z)/z^2$、$\sin z/z^3$、$\sin(1/z)$。
 \item 设 $f$ 在去心邻域全纯，且 $(z-a)f(z)\to0$，证明 $a$ 为可去奇点。
 \item 判断 $z^3+1$、$z/(z^2+1)$、$\sin z$ 在无穷远点的类型。
 \item 解释 $\operatorname{Log}z$ 的原点为何不属于本章的孤立奇点分类，并与 $e^{1/z}$ 比较。
\end{enumerate}

